%% Springer Nature sn-jnl format for Artificial Intelligence and Law
%% Target: subscription model (no APC)
%% Review: double-blind
%%
\documentclass[pdflatex,sn-basic,Numbered]{sn-jnl}

\usepackage{multirow}
\usepackage{amsmath,amssymb,amsfonts}
\usepackage{booktabs}
\usepackage{xcolor}

\raggedbottom

\begin{document}

\title[Neuro-Symbolic Legal Reasoning for Indonesian Legal Pluralism]{A Neuro-Symbolic Methodology for Legal Reasoning in Indonesian Legal Pluralism: Expert-Verified Customary Law Rules and Pilot Evaluation}

%%=============================================================%%
%% Double-blind: author info hidden during review              %%
%%=============================================================%%
\author*[1]{\fnm{Anonymous} \sur{Author 1}}\email{anonymous@example.com}
\author[1]{\fnm{Anonymous} \sur{Author 2}}

\affil*[1]{\orgdiv{[Withheld]}, \orgname{[Withheld]}, \orgaddress{\city{[Withheld]}, \country{[Withheld]}}}

\abstract{We present a neuro-symbolic methodology for legal reasoning in Indonesian legal pluralism, where codified national law coexists with multiple customary (\textit{adat}) normative systems that can conflict.
Our primary contribution is methodological and resource-oriented: (1) an expert-verified customary-law rule base encoded in Answer Set Programming (ASP) across three domains---Minangkabau (25 rules), Bali (34 rules), Jawa (36 rules)---together with (2) an auditable benchmark protocol with manifest-based provenance tracking and automated claim gating, and (3) a transparent pilot evaluation framework that documents both positive and negative findings.
An optional LLM adjudication layer serves as a decision interface over symbolic outputs.
On 74 dual-expert-labeled scenarios (70 evaluable, 4 disputed), we compare ASP-only (58.6\%, Wilson 95\% CI [0.469, 0.694]), ASP+local-LLM (64.3\%, CI [0.526, 0.745]), and ASP+API-LLM (68.6\%, CI [0.570, 0.782]).
Observed gains are directionally positive but statistically inconclusive at pilot scale (all McNemar $p \geq 0.167$, $n=70$).
Inter-system agreement is substantial (Fleiss' $\kappa = 0.638$), consistent with symbolic-layer anchoring across backends.
We document key negative results---including a rule over-specification paradox where expanding rule coverage \textit{decreases} accuracy, backend-dependent performance sensitivity, and complete failure on uncertainty detection---and discuss their methodological implications for neuro-symbolic legal AI.
This work is positioned as a benchmark-and-methodology contribution; reaching statistical power $\geq$0.8 requires approximately 344 evaluable cases.}

\keywords{Legal pluralism \and Neuro-symbolic AI \and Answer Set Programming \and Customary law \and Indonesian law \and Legal NLP \and Pilot benchmark}

\maketitle

\section{Introduction}\label{sec:intro}

Indonesia poses a challenging test bed for legal AI because legal decisions often involve legal pluralism: codified national law coexists with multiple customary (\textit{adat}) normative systems.
In such settings, legal reasoning requires more than retrieving relevant text.
The system must resolve conflicts across regimes, handle exceptions, and justify why one normative source is prioritized in a given case.

\subsection{Cross-disciplinary context}\label{subsec:crossdisc}

This paper is intentionally written for three communities that rarely share technical language: legal scholars, customary-law and culture scholars, and computer scientists.
We therefore clarify core terms early.
\textit{Natural Language Processing (NLP)} refers to computational methods for processing legal text \cite{guha2023legalbench}; \textit{Large Language Models (LLMs)} are neural models used here for text-level adjudication \cite{brown2020language}; and \textit{Answer Set Programming (ASP)} is a symbolic logic formalism for explicit rule-based reasoning under defaults and exceptions \cite{lifschitz2019asp,gebser2014clingo}.
Our \textit{neuro-symbolic} design combines these roles: neural components map text into usable representations, while symbolic components provide auditable legal constraints \cite{garcez2023neurosymbolic}.

\subsection{Motivation for each audience}\label{subsec:motivation}

For legal readers, the core issue is doctrinal defensibility in plural-law conflicts, where opaque end-to-end neural predictions are difficult to audit.
For customary-law and cultural readers, the core issue is representation fidelity---\textit{adat} norms cannot be treated as informal noise if they are normatively operative in practice \cite{hooker1978adat,burns2004leiden}.
For computer-science readers, the core issue is methodological---current legal NLP benchmarks mostly assume one dominant legal regime, which underestimates conflict-handling complexity in plural systems.

\subsection{Research gap and contributions}\label{subsec:gap}

Prior legal NLP benchmarks focus mostly on monolithic legal systems and text-centric tasks such as judgment prediction, extraction, and summarization.
Prior Indonesian legal-pluralism scholarship is rich in doctrinal and anthropological analysis but is typically not encoded as executable reasoning artifacts.
The missing bridge is an auditable benchmark line that jointly provides: (1) expert-verified Indonesian \textit{adat} rules, (2) formal executable encoding, and (3) transparent pilot evaluation under explicit governance constraints.

This paper addresses that bridge as a methodology-first neuro-symbolic study.
We use ASP as the rule-centric reasoning core and treat the LLM layer as an adjudication interface.
Our concrete contributions are:
\begin{enumerate}
  \item A formally encoded, expert-verified ASP rule base for three Indonesian customary-law domains: Minangkabau (25 rules), Bali (34 rules), and Jawa (36 rules). Of 95 verified rules, 71 are currently active in the reasoning engine; the remaining 24 caused accuracy regression when added (Section~\ref{subsec:overspec}).
  \item A pilot benchmark resource and auditable reporting protocol over 74 dual-expert-labeled scenarios (70 evaluable, 4 disputed), including manifest-based provenance tracking and automated claim gating.
  \item Transparent documentation of negative results, including the rule over-specification paradox, backend sensitivity, complete D-label failure, and pilot-scale statistical limits.
  \item Cross-backend pilot evidence showing directionally positive but statistically inconclusive gains, plus substantial inter-system agreement (Fleiss' $\kappa = 0.638$), consistent with symbolic-layer anchoring.
\end{enumerate}

\subsection{Scope and paper roadmap}\label{subsec:scope}

This manuscript focuses on rule formalization, expert verification workflow, and pilot-scale evaluation.
The current results are explicitly preliminary: statistical power is limited ($n=70$ evaluable; all McNemar $p \geq 0.167$), and the same 70 evaluable cases have been used during development and reporting.
We frame the paper as a benchmark-and-methodology contribution, deferring strong efficacy claims to future unseen-case evaluation.

Readers interested in legal and cultural framing may focus on Sections~\ref{sec:related}, \ref{sec:system}, \ref{sec:dataset}, and \ref{sec:limitations}.
Readers focused on computational design and evaluation may focus on Sections~\ref{sec:system}, \ref{sec:metrics}, and \ref{sec:results}.
Section~\ref{sec:discussion} connects findings across both audiences.

\section{Related Work}\label{sec:related}

This study sits at the intersection of legal pluralism scholarship, legal NLP, and neuro-symbolic AI.
Each body of work solves part of the pipeline needed for auditable plural-law reasoning, but none addresses the full chain.

\subsection{Legal pluralism and customary-law scholarship}\label{subsec:pluralism}

Indonesian legal anthropology and doctrinal studies have long documented that national law and \textit{adat} norms co-govern many real disputes \cite{hooker1978adat,burns2004leiden}.
This literature explains normative coexistence, local legitimacy, and context-dependent conflict handling.
However, the knowledge remains non-executable: rich for doctrinal interpretation but not expressible as machine-checkable reasoning rules.
Our work bridges this gap by encoding expert-verified \textit{adat} norms as formal ASP rules while preserving their domain-specific semantics.

\subsection{Legal NLP benchmarks}\label{subsec:legalnlp}

Mainstream legal NLP has advanced tasks such as summarization \cite{sie-etal-2024-summarizing,tyss-etal-2024-lexsumm}, judgment prediction \cite{nigam-etal-2024-rethinking,liu-etal-2024-enhancing-legal,xie-etal-2024-clc}, contract understanding \cite{mali-etal-2024-information,narendra-etal-2024-enhancing,kwak-etal-2024-classify}, and domain-specific retrieval \cite{hou-etal-2024-gaps,taranukhin-etal-2024-empowering}.
Collaborative benchmarks like LegalBench \cite{guha2023legalbench} and domain-adapted models such as SaulLM \cite{colombo2024saullm} and DeBERTa fine-tunes \cite{tran-etal-2024-deberta} have raised the baseline for monolithic-regime tasks.
However, these datasets and models assume a single dominant authority structure and therefore do not evaluate cross-regime conflict resolution---the core challenge in plural-law settings.
Our benchmark specifically targets this gap with a four-label taxonomy (national, customary, synthesis, unclear) that requires the system to reason about normative conflicts.

\subsection{LLM reasoning and retrieval}\label{subsec:llm}

Advances in prompting strategies---chain-of-thought \cite{wei2022cot}, self-consistency \cite{wang2022selfconsistency}, ReAct \cite{yao2023react}, tree-of-thought \cite{yao2023tree}, and Reflexion \cite{shinn2023reflexion}---have improved LLM reasoning.
Retrieval-augmented generation extends grounding by injecting external context \cite{lewis2020rag,asai2023selfrag,gao2023ragsurvey,sarthi2024raptor,he2024gretriever}.
Yet in plural-law settings, retrieval and fluent generation are insufficient: the system must also enforce explicit conflict constraints and exception logic that these approaches do not natively provide.

\subsection{Neuro-symbolic AI for legal reasoning}\label{subsec:neurosym}

Neuro-symbolic AI combines neural flexibility with symbolic control \cite{garcez2023neurosymbolic}.
ASP, in particular, provides non-monotonic rule semantics that align with legal defaults, exceptions, and contradictions \cite{lifschitz2019asp,gebser2014clingo}.
This is attractive for legal stakeholders because reasoning traces are inspectable and norm references auditable.
However, existing neuro-symbolic work in legal AI has not addressed multi-regime conflict resolution with expert-verified customary-law rules.
Our system fills this gap by using ASP as a \textit{shared reasoning anchor} across multiple LLM backends, with the pilot evidence showing substantial inter-system agreement (Fleiss' $\kappa = 0.638$).

\section{System Overview}\label{sec:system}

The system is rule-centric: ASP formalization of expert-verified customary law is the primary contribution.
The LLM-based adjudication layer is a practical decision interface, not the central novelty.

This design is motivated by three requirements specific to legal reasoning in pluralistic jurisdictions: (1) \textit{auditability}---ASP rules provide explicit, inspectable reasoning traces; (2) \textit{non-monotonic reasoning}---ASP's native support for defaults and exceptions accommodates defeasible customary norms; and (3) \textit{deployment flexibility}---the system operates without API calls in offline mode using deterministic heuristics.

\subsection{Pipeline architecture}\label{subsec:arch}

Given a legal scenario $x$ in natural language, the system predicts one of four policy labels: \textbf{A} (mainly national-law resolution), \textbf{B} (mainly customary-law resolution), \textbf{C} (synthesis of national + customary law), or \textbf{D} (clarification required).
The pipeline proceeds in three stages:

\begin{enumerate}
  \item \textbf{Keyword router and fact extraction.} A keyword-based classifier identifies the customary-law domain (Minangkabau, Bali, Jawa, or national) and extracts legal predicates from the scenario text. For example, the presence of terms such as \textit{pusako tinggi} (high ancestral property) and \textit{matrilineal} triggers the Minangkabau domain and generates ASP facts such as \texttt{jenis\_harta(pusako\_tinggi)} and \texttt{garis\_keturunan(matrilineal)}.
  \item \textbf{ASP rule engine (Clingo).} The extracted facts are passed to the Clingo solver together with domain-specific ASP rules. The engine performs non-monotonic inference to detect regime conflicts, apply defaults, and propose a label. If both national and customary predicates are active and trigger contradictory conclusions, the engine signals \texttt{konflik\_terdeteksi(ya)}.
  \item \textbf{LLM adjudication (optional).} In ASP+LLM mode, the symbolic output---including the proposed label, conflict signals, and active rules---is passed to an LLM for final adjudication. In ASP-only mode, a deterministic heuristic supervisor makes the final decision without any LLM call.
\end{enumerate}

Figure~\ref{fig:architecture} illustrates this pipeline.

\begin{figure}[t]
\centering
\small
\begin{verbatim}
 +------------------+     +------------------+     +------------------+
 |   Legal Scenario |---->| Keyword Router + |---->|  ASP Rule Engine |
 |   (natural lang) |     | Fact Extraction  |     |    (Clingo)      |
 +------------------+     +------------------+     +------------------+
                                                          |
                          +-------------------------------+
                          |                               |
                          v                               v
                 +------------------+           +------------------+
                 | LLM Adjudication |           |  Offline Heurist.|
                 |  (ASP+LLM mode) |           |  (ASP-only mode) |
                 +------------------+           +------------------+
                          |                               |
                          +-------------------------------+
                                      |
                                      v
                           +--------------------+
                           | Final Label (A/B/C/D)
                           | + Reasoning Trace  |
                           +--------------------+
\end{verbatim}
\caption{Pipeline architecture. The ASP rule engine acts as the shared reasoning core. The LLM adjudication layer is optional; in ASP-only mode, a deterministic heuristic produces the final label.}\label{fig:architecture}
\end{figure}

\subsection{ASP rule design}\label{subsec:asp-rules}

Each domain-specific rule file encodes expert-verified customary-law norms as ASP rules.
Rules follow a consistent pattern: domain facts trigger inheritance classifications, and conflict detection rules fire when national-law and customary-law conclusions co-occur.
Listing~\ref{lst:asp-example} shows a simplified Minangkabau rule for matrilineal high-property inheritance.

\begin{figure}[t]
\centering
\small
\begin{verbatim}
%% Minangkabau: high ancestral property (pusako tinggi)
%% follows matrilineal line exclusively
hak_waris(perempuan_suku) :-
    domain(minangkabau),
    jenis_harta(pusako_tinggi),
    garis_keturunan(matrilineal).

%% Conflict detection: national civil code vs adat
konflik_terdeteksi(ya) :-
    hak_waris_nasional(X),
    hak_waris(Y),
    X != Y.

%% Default label proposal
label_proposal(b) :-
    domain(minangkabau),
    hak_waris(_),
    not konflik_terdeteksi(ya).

label_proposal(c) :-
    konflik_terdeteksi(ya).
\end{verbatim}
\caption{Simplified ASP rules for Minangkabau matrilineal inheritance. \texttt{hak\_waris} assigns inheritance rights; \texttt{konflik\_terdeteksi} detects conflicts between national and customary conclusions.}\label{lst:asp-example}
\end{figure}

\subsection{Expert-verified rule base}\label{subsec:rulebase}

Table~\ref{tab:rulebase} summarizes the expert-verified customary-law rules across three domains.
Of 95 total verified rules, 71 are currently encoded in ASP.
The remaining 24 were rolled back after empirical testing showed they caused accuracy regression (Section~\ref{subsec:overspec}).

\begin{table}[t]
\caption{Expert-verified customary-law rule base.}\label{tab:rulebase}
\begin{tabular*}{\textwidth}{@{\extracolsep\fill}lrrl}
\toprule
Domain & Verified & Active in ASP & Scope \\
\midrule
Minangkabau & 25 & 17 & Matrilineal inheritance and \textit{pusako} norms \\
Bali & 34 & 30 & \textit{Purusa/sentana} and \textit{druwe} classifications \\
Jawa & 36 & 24 & Bilateral inheritance and \textit{gono-gini} patterns \\
\midrule
Total & 95 & 71 & Expert-verified customary law base \\
\botrule
\end{tabular*}
\end{table}

\section{Dataset and Expert Protocol}\label{sec:dataset}

\subsection{Benchmark construction}\label{subsec:benchmark}

The benchmark dataset was constructed in two phases: an initial batch of 24 cases and an expanded batch of 50 new cases.
Each case was independently labeled by two qualified legal experts (identities withheld for double-blind review).
Table~\ref{tab:status} summarizes the dataset status, and Table~\ref{tab:labeldist} shows the label distribution.

\begin{table}[t]
\caption{Operational dataset status.}\label{tab:status}
\begin{tabular*}{\textwidth}{@{\extracolsep\fill}ll}
\toprule
Item & Value \\
\midrule
Total benchmark cases & 74 \\
Evaluable (agreed) cases & 70 \\
Disputed cases (excluded from accuracy) & 4 \\
Initial batch agreement (24 cases) & 58.3\% (Cohen's $\kappa = 0.394$) \\
Expanded batch agreement (50 cases) & 94.0\% (47/50) \\
\botrule
\end{tabular*}
\end{table}

\begin{table}[t]
\caption{Label distribution in the evaluable benchmark subset ($N=70$).}\label{tab:labeldist}
\begin{tabular*}{\textwidth}{@{\extracolsep\fill}lrr}
\toprule
Gold Label & Count & Percentage \\
\midrule
A (National) & 6 & 8.6\% \\
B (Customary) & 31 & 44.3\% \\
C (Synthesis) & 31 & 44.3\% \\
D (Unclear) & 2 & 2.9\% \\
\botrule
\end{tabular*}
\end{table}

\subsection{Rubric refinement}\label{subsec:rubric}

Initial inter-rater agreement in the 24-case batch was moderate ($\kappa = 0.394$; 14/24 agreement), with recurring ambiguity at the A/B/C decision boundary.
The annotation protocol was tightened in stages: (i) locked A/B/C/D label definitions, (ii) explicit boundary tests for A vs.\ C, and (iii) a mandatory decision checklist (dominance, dual-requirement, missing-fact check).
The expanded 50-case packet applied this refined rubric on a disjoint case pool and explicitly instructed raters not to inspect prior labels.
Agreement on the expanded pool reached 47/50 (94.0\%), with three remaining disagreements at the B/C boundary marked as disputed.
A complete audit trail of this refinement process is maintained in the project methodology log.

\subsection{Adjudication policy}\label{subsec:adjudication}

Only agreed cases are used for quantitative accuracy reporting.
Disputed cases are excluded and queued for third-expert tiebreak.
The benchmark manifest is treated as the reporting authority for snapshot counts.

\section{Metrics}\label{sec:metrics}

We explicitly define pilot metrics to avoid ambiguous reporting.

\subsection{Agreement coverage}\label{subsec:consensus}

Let $N$ be total active cases and $N_a$ the agreed/evaluable subset:
\begin{equation}
\mathrm{ConsensusStrength} = \frac{N_a}{N}
\end{equation}
For the current snapshot, $(N_a, N) = (70, 74)$, so $\mathrm{ConsensusStrength} = 0.946$.

\subsection{Wilson confidence interval}\label{subsec:wilson}

For observed accuracy $\hat{p}$ with sample size $n$ and $z=1.96$ (95\% CI):
\begin{equation}
\frac{\hat{p} + \frac{z^2}{2n} \pm z\sqrt{\frac{\hat{p}(1-\hat{p})}{n}+\frac{z^2}{4n^2}}}{1+\frac{z^2}{n}}
\label{eq:wilson}
\end{equation}
Wilson intervals are preferred over normal approximation for small pilot datasets where $\hat{p}$ may be near the boundary \cite{wilson1927probable}.

\section{Experiments and Results}\label{sec:results}

\subsection{Ablation: ASP-only vs.\ ASP+LLM}\label{subsec:ablation}

We compare three operating modes on the 70 evaluable cases:
\begin{itemize}
  \item \textbf{ASP-only}: keyword router $\rightarrow$ ASP rule engine $\rightarrow$ offline heuristic supervisor (no LLM calls).
  \item \textbf{ASP+local-LLM}: as above, with LLM adjudication via a local 7B-parameter model (deepseek-r1 via Ollama, \texttt{temperature=0}).
  \item \textbf{ASP+API-LLM}: as above, with LLM adjudication via a commercial API (DeepSeek-Chat, \texttt{temperature=0}).
\end{itemize}

\begin{table}[t]
\caption{Ablation comparison on the 70-case evaluable benchmark (Table~\ref{tab:ablation}).}\label{tab:ablation}
\begin{tabular*}{\textwidth}{@{\extracolsep\fill}lccc}
\toprule
Mode & Accuracy & Wilson 95\% CI & Correct / Total \\
\midrule
ASP-only & 58.6\% & [0.469, 0.694] & 41 / 70 \\
ASP+local-LLM & 64.3\% & [0.526, 0.745] & 45 / 70 \\
ASP+API-LLM & \textbf{68.6\%} & [0.570, 0.782] & 48 / 70 \\
\botrule
\end{tabular*}
\end{table}

The LLM layer shows an observed improvement of +5.7 pp (local) and +10.0 pp (API) over ASP-only (Table~\ref{tab:ablation}).
However, all pairwise McNemar tests are non-significant ($p \geq 0.167$, $n=70$; see Table~\ref{tab:mcnemar}); these differences should be interpreted as pilot observations pending larger-scale evaluation.

\subsection{Rule over-specification paradox}\label{subsec:overspec}

A central negative finding is that increasing ASP rule coverage from 71 to 95 rules (100\% of expert-verified rules) paradoxically \textit{decreased} accuracy.
Table~\ref{tab:overspec} summarizes the effect.

\begin{table}[t]
\caption{Rule over-specification experiment: accuracy before and after activating 24 additional ASP rules.}\label{tab:overspec}
\begin{tabular*}{\textwidth}{@{\extracolsep\fill}lccc}
\toprule
Configuration & Active Rules & Accuracy & $\Delta$ vs.\ 71 rules \\
\midrule
ASP+local-LLM (71 rules) & 71 & 64.3\% & --- \\
ASP+local-LLM (95 rules) & 95 & 57.1\% & $-$7.1 pp \\
\botrule
\end{tabular*}
\footnotetext{The 24 additional rules were rolled back after this result. All canonical results use 71 active rules.}
\end{table}

Analysis revealed that the additional rules were overly general, introducing spurious \textit{adat}-dominance signals that the LLM amplified rather than corrected.
Conflict cases (gold=C) were systematically misclassified as B (pure customary).
This suggests that neuro-symbolic integration requires calibrating symbolic information exposure rather than maximizing rule coverage---a finding we term the \textit{rule over-specification paradox}.

\subsection{Per-label performance}\label{subsec:perlabel}

\begin{table}[t]
\caption{Per-label precision (P), recall (R), and F1 on the 70-case benchmark.}\label{tab:labelperf}
\begin{tabular*}{\textwidth}{@{\extracolsep\fill}lccccccccc}
\toprule
 & \multicolumn{4}{@{}c@{}}{ASP-only} & \multicolumn{4}{@{}c@{}}{ASP+API-LLM} \\
\cmidrule{2-5}\cmidrule{6-9}
Label & P & R & F1 & Supp. & P & R & F1 & Supp. \\
\midrule
A & 0.36 & 0.67 & 0.47 & 6 & 0.40 & 0.67 & 0.50 & 6 \\
B & 0.66 & 0.61 & 0.63 & 31 & 0.69 & 0.77 & 0.73 & 31 \\
C & 0.62 & 0.58 & 0.60 & 31 & 0.80 & 0.65 & 0.71 & 31 \\
D & 0.00 & 0.00 & 0.00 & 2 & 0.00 & 0.00 & 0.00 & 2 \\
\botrule
\end{tabular*}
\end{table}

\subsection{Error analysis: the 13 hard cases}\label{subsec:hardcases}

We analyze 13 cases (18.6\%) where all three system variants fail simultaneously (Table~\ref{tab:hardcases-dist}).

\begin{table}[t]
\caption{Label distribution in 13 hard cases versus overall benchmark ($N=70$).}\label{tab:hardcases-dist}
\begin{tabular*}{\textwidth}{@{\extracolsep\fill}lccc}
\toprule
Gold Label & Hard Cases & Overall & Ratio \\
\midrule
C (Conflict) & 6 (46.2\%) & 31 (44.3\%) & 1.04$\times$ \\
B (Adat) & 5 (38.5\%) & 31 (44.3\%) & 0.87$\times$ \\
A (National) & 1 (7.7\%) & 6 (8.6\%) & 0.90$\times$ \\
D (Unclear) & 1 (7.7\%) & 2 (2.9\%) & 2.67$\times$ \\
\botrule
\end{tabular*}
\end{table}

\begin{table}[t]
\caption{Failure patterns in 13 hard cases where all three systems fail.}\label{tab:failure-patterns}
\begin{tabular*}{\textwidth}{@{\extracolsep\fill}llcc}
\toprule
Pattern & Description & Count & Pct \\
\midrule
C$\rightarrow$B & Conflict misclassified as Adat & 4 & 30.8\% \\
B$\rightarrow$A & Adat misclassified as National & 3 & 23.1\% \\
B$\rightarrow$C & Adat misclassified as Conflict & 2 & 15.4\% \\
C$\rightarrow$A & Conflict misclassified as National & 2 & 15.4\% \\
D$\rightarrow$C & Unclear misclassified as Conflict & 1 & 7.7\% \\
A$\rightarrow$C & National misclassified as Conflict & 1 & 7.7\% \\
\botrule
\end{tabular*}
\end{table}

The dominant pattern is C$\rightarrow$B misclassification (30.8\% of hard cases): when customary-law keywords dominate the scenario text, the keyword router overwhelms the conflict signal (Table~\ref{tab:failure-patterns}).
A two-layer error audit across all 29 C$\rightarrow$B events (pooled across three configurations) reveals that 23 of 29 (79.3\%) occur when the ASP engine produces no conflict signal---indicating upstream router/fact-extraction failure---while 6 of 29 (20.7\%) occur despite a conflict signal, indicating adjudicator-level failure.

\subsection{Per-domain analysis}\label{subsec:domain}

\begin{table}[t]
\caption{Per-domain accuracy across three system configurations.}\label{tab:domain}
\begin{tabular*}{\textwidth}{@{\extracolsep\fill}lrrrr}
\toprule
Domain & $N$ & ASP-only & ASP+local & ASP+API \\
\midrule
Minangkabau & 21 & 71.4\% & 76.2\% & 71.4\% \\
Bali & 21 & 71.4\% & 81.0\% & 76.2\% \\
Jawa & 17 & 35.3\% & 29.4\% & 52.9\% \\
Nasional & 7 & 42.9\% & 71.4\% & 71.4\% \\
Lintas & 4 & 50.0\% & 50.0\% & 50.0\% \\
\midrule
Overall & 70 & 58.6\% & 64.3\% & 68.6\% \\
\botrule
\end{tabular*}
\end{table}

The Jawa domain shows the weakest ASP-only performance (35.3\%) but the largest LLM delta (+17.6 pp with ASP+API-LLM; Table~\ref{tab:domain}).
This reveals that the LLM is \textit{compensating} for weak symbolic rules rather than augmenting strong ones---an important finding discussed further in Section~\ref{sec:discussion}.
Conversely, ASP+local-LLM \textit{decreases} Jawa accuracy to 29.4\% ($-$5.9 pp), indicating that LLM capability differences interact with domain-specific rule quality.

\subsection{Cross-backend validation}\label{subsec:crossval}

\begin{table}[t]
\caption{Cross-validation across three backends ($N=70$). All runs use identical rule state (71 rules) and \texttt{temperature=0}.}\label{tab:crossval}
\begin{tabular*}{\textwidth}{@{\extracolsep\fill}lcccc}
\toprule
Mode & Accuracy & Wilson 95\% CI & Cohen's $\kappa$ & Agreement \\
\midrule
ASP-only & 58.6\% & [0.469, 0.694] & 0.331 & Fair \\
ASP+local-LLM & 64.3\% & [0.526, 0.745] & 0.418 & Moderate \\
ASP+API-LLM & \textbf{68.6\%} & [0.570, 0.782] & 0.483 & Moderate \\
\botrule
\end{tabular*}
\end{table}

\begin{table}[t]
\caption{McNemar exact binomial test results ($N=70$).}\label{tab:mcnemar}
\begin{tabular*}{\textwidth}{@{\extracolsep\fill}lccc}
\toprule
Comparison & $p$-value & Discordant & Favored \\
\midrule
ASP-only vs.\ local-LLM & $p = 0.344$ & 10 & local (+4) \\
ASP-only vs.\ API-LLM & $p = 0.167$ & 19 & API (+7) \\
local-LLM vs.\ API-LLM & $p = 0.549$ & 11 & API (+3) \\
\botrule
\end{tabular*}
\end{table}

\textbf{No differences reach statistical significance} at $\alpha = 0.05$ (Table~\ref{tab:mcnemar}).

\subsection{Inter-system agreement}\label{subsec:agreement}

\begin{table}[t]
\caption{Cross-model agreement among three configurations ($N=70$).}\label{tab:agreement}
\begin{tabular*}{\textwidth}{@{\extracolsep\fill}lcc}
\toprule
Pattern & Count & Pct \\
\midrule
All 3 agree (correct) & 35 & 50.0\% \\
All 3 agree (wrong) & 13 & 18.6\% \\
\midrule
\textbf{All 3 agree (total)} & \textbf{48} & \textbf{68.6\%} \\
2 of 3 agree, 1 different & 19 & 27.1\% \\
All 3 different & 3 & 4.3\% \\
\botrule
\end{tabular*}
\end{table}

Fleiss' $\kappa = 0.638$ (substantial agreement) with only 4.3\% complete disagreement (Table~\ref{tab:agreement}), consistent with the ASP rule layer acting as a strong prior across backends.
Majority vote on the 19 partial-agreement cases achieves only 47.4\% accuracy (9/19), indicating that when systems disagree, the majority is no more likely correct than chance.

\subsection{Exploratory backend sensitivity}\label{subsec:exploratory}

Two additional local backends tested after the canonical freeze show divergent behavior under the same rule state and prompts.
A 20B-parameter open-source model matched the canonical local result (45/70, 64.3\%), while a 14B-parameter model dropped to 38/70 (54.3\%), performing \textit{worse} than ASP-only.
Both share a systematic B$\rightarrow$A bias pattern (6 identical failures) where national-law signals in ASP traces are over-interpreted as national-law dominance.
All paired McNemar comparisons remain non-significant (20B: $p = 0.646$; 14B: $p = 0.480$).

\section{Discussion}\label{sec:discussion}

\subsection{What the methodology reveals}\label{subsec:disc-method}

The primary contribution of this work is not an accuracy number but a \textit{methodology} for building auditable neuro-symbolic legal reasoning in plural-law settings.
Three methodological insights emerge from the pilot.

First, the rule over-specification paradox (Section~\ref{subsec:overspec}) demonstrates that maximizing symbolic coverage can \textit{harm} neuro-symbolic performance.
When 24 additional ASP rules were activated, accuracy dropped by 7.1 pp (Table~\ref{tab:overspec}) because overly general rules introduced spurious \textit{adat}-dominance signals that the LLM amplified rather than corrected.
This suggests that rule base design for neuro-symbolic systems should prioritize \textit{precision over coverage}---a principle not widely discussed in the neuro-symbolic literature.

Second, the per-domain analysis (Section~\ref{subsec:domain}) reveals an asymmetry: the LLM layer provides the largest gains where symbolic rules are weakest (Jawa: +17.6 pp with API-LLM), but contributes minimally or even negatively where rules are already strong (Minangkabau: $-$4.8 pp with API-LLM vs.\ ASP+local).
This is consistent with a \textit{compensation} dynamic rather than \textit{augmentation}: the neural layer substitutes for missing symbolic knowledge rather than enhancing existing symbolic reasoning.
For system designers, this implies that investment in rule quality for weak domains may yield more benefit than upgrading LLM backends.

Third, the rubric refinement arc---from 58.3\% inter-rater agreement ($\kappa = 0.394$) to 94.0\% (47/50)---illustrates that legal pluralism classification tasks are highly sensitive to boundary definitions.
The improvement was achieved through locked label definitions, explicit boundary tests, and decision checklists, not through changing the underlying legal substance.
This finding is relevant beyond our specific system: any legal AI benchmark for plural-law settings must invest significantly in rubric engineering.

\subsection{Connecting the error threads}\label{subsec:disc-errors}

The error analysis reveals three interconnected failure modes.
The dominant C$\rightarrow$B misclassification stems primarily from upstream router failures (79.3\% of cases): when ethnic keywords dominate, the keyword-based router misses implicit conflict signals.
The B$\rightarrow$A bias observed in open-source backends represents a complementary failure: these models over-interpret national-law terminology in ASP traces.
The complete D-label failure (0\% recall, $n=2$) indicates the absence of an abstention mechanism.

These three patterns converge on a single design principle: \textit{the weakest link in a neuro-symbolic pipeline is the text-to-symbol interface}, not the symbolic reasoning or the neural adjudication in isolation.
Router improvements and abstention mechanisms are therefore higher-priority targets than rule expansion or LLM upgrades.

\subsection{Statistical power and pilot framing}\label{subsec:disc-power}

All pairwise McNemar tests are non-significant ($p \geq 0.167$), and estimated statistical power is approximately 0.3 for detecting a 10 pp difference at $n=70$.
Reaching power $\geq 0.8$ would require approximately 344 evaluable cases, following standard sample size estimation for McNemar's test \cite{lachin1992power}.
Non-significance does not mean the observed differences are zero---it means they cannot be distinguished from zero at current scale.
The substantial inter-system agreement (Fleiss' $\kappa = 0.638$) provides complementary evidence that system behavior is reproducible across backends, even though accuracy differences are inconclusive.

\subsection{Implications for neuro-symbolic legal AI}\label{subsec:disc-implications}

For the neuro-symbolic AI community, our negative results carry methodological weight.
The rule over-specification paradox challenges the intuition that more rules always help: in neuro-symbolic integration, symbolic outputs serve as \textit{context} for neural decision-making, and noisy context degrades neural performance.
The backend sensitivity finding---where one local model improves over ASP-only while another degrades it---demonstrates that the neural component is not interchangeable, and that integration testing must be backend-specific.

For the legal AI community, the rubric refinement process demonstrates that inter-rater agreement in plural-law settings can be improved from 58.3\% to 94.0\% through structured boundary definitions, suggesting that apparent disagreement often reflects rubric ambiguity rather than genuine legal indeterminacy.

\section{Limitations}\label{sec:limitations}

\subsection{Statistical power}\label{subsec:lim-power}

The primary limitation is $n=70$ evaluable cases with estimated power $\approx 0.3$ for a 10 pp effect.
All findings should be treated as pilot observations.

\subsection{Development contamination}\label{subsec:lim-contam}

The 70 evaluable cases were used for both prompt tuning (accuracy rose from 54\% to 68.6\% during development) and final evaluation.
A held-out test set is defined by protocol but awaits new expert annotations.

\subsection{Rule parity gaps}\label{subsec:lim-parity}

The ASP rule base captures only a subset of expert-verified norms: Minangkabau 5/25 fully covered, Bali 21/34, Jawa 15/36.
Procedural and contemporary norms are under-represented, particularly in the Jawa domain (35.3\% ASP-only accuracy).

\subsection{Uncertainty detection}\label{subsec:lim-dlabel}

D-label recall is 0\% across all configurations ($n=2$).
The system lacks an explicit abstention mechanism.

\subsection{Scope and generalizability}\label{subsec:lim-scope}

All cases are in Indonesian across three specific customary systems.
Generalization to other legal pluralism contexts is not established.

\section{Conclusion}\label{sec:conclusion}

This paper contributes a methodology and pilot benchmark for neuro-symbolic legal reasoning in Indonesian legal pluralism: an expert-verified customary-law rule base encoded in ASP, an auditable evaluation protocol with manifest-based provenance, and transparent reporting of both positive and negative findings.

On the 70-case evaluable benchmark, ASP+LLM configurations show observed accuracies of 64.3\%--68.6\% versus 58.6\% for ASP-only (directional gain: +5.7 to +10.0 pp), with all McNemar tests non-significant ($n=70$, $p \geq 0.167$).
Substantial inter-system agreement (Fleiss' $\kappa = 0.638$) is consistent with symbolic-layer anchoring.

Key negative findings---the rule over-specification paradox, backend-dependent sensitivity, and complete D-label failure---provide methodological insights for future neuro-symbolic legal AI systems.
The dominant error source is the text-to-symbol interface rather than the symbolic or neural components in isolation.

Future work should: (1) expand the benchmark to $\geq$344 cases for adequate statistical power; (2) establish a held-out test set from new expert annotations; (3) redesign the router for implicit conflict detection and D-label abstention; and (4) investigate the compensation-vs-augmentation dynamic across domains with varying rule quality.

\backmatter

\bmhead{Acknowledgments}

[Withheld for double-blind review.]

\section*{Declarations}

\begin{itemize}
\item \textbf{Funding}: [Withheld for double-blind review.]
\item \textbf{Conflict of interest}: The authors declare no competing interests.
\item \textbf{Ethics approval}: Not applicable (no human subjects; expert annotators participated voluntarily as research collaborators).
\item \textbf{Consent to participate}: Not applicable.
\item \textbf{Consent for publication}: Not applicable.
\item \textbf{Data availability}: The dataset and benchmark artifacts will be made available in a public repository upon acceptance.
\item \textbf{Code availability}: Source code will be made available in a public repository upon acceptance.
\item \textbf{Authors' contributions}: [Withheld for double-blind review.]
\end{itemize}

\bibliography{references}

\end{document}
